\begin{scanpage}{0081}{64}
\[
\begin{aligned}
E_\vartheta=H_\varphi={}&\frac{ae^{i\pi/4}\cos\vartheta}{\sqrt{2\pi ka\sqrt{\lambda^2+\mu^2}}}\frac{e^{ikR}}R\biggl\{-H_{0x_1}\left[\frac{\cos(\psi_1-\varphi)}{\sin\alpha_1\cos\varphi_1}\right.\\
&\left.+\frac{\sin(\psi_1-\varphi)}{\sin^2\alpha_1}\cos\alpha_1\operatorname{tg}\varphi_1\right][g^1(\varphi_1,\varphi_1^0)e^{-ika\sqrt{\lambda^2+\mu^2}}\\
&-ig^1(\varphi_2,\varphi_2^0)e^{ika\sqrt{\lambda^2+\mu^2}}]+E_{0x_1}\frac{\sin(\psi_1-\varphi)}{\sin^2\alpha_1}\times\\
&\times[f^1(\varphi_1,\varphi_1^0)e^{-ika\sqrt{\lambda^2+\mu^2}}-if^1(\varphi_2,\varphi_2^0)e^{ika\sqrt{\lambda^2+\mu^2}}]\biggr\}.
\end{aligned}\tag{11.31}
\]
The resulting expressions are valid when $ka\sqrt{\lambda^2+\mu^2}\gg1$. They may be slightly simplified to
\[
\begin{aligned}
E_\varphi=-H_\vartheta={}&\frac{ae^{i\pi/4}}{\sqrt{2\pi ka\sqrt{\lambda^2+\mu^2}}}\frac{e^{ikR}}R\times\\
&\times\biggl\{-H_{0x_1}\cos\vartheta\frac{\sin(\psi_1-\varphi)}{\sin^2\alpha_1}[g^1(\varphi_1,\varphi_1^0)e^{-ika\sqrt{\lambda^2+\mu^2}}\\
&-ig^1(\varphi_2,\varphi_2^0)e^{ika\sqrt{\lambda^2+\mu^2}}]-E_{0x_1}\frac{\cos(\psi_1-\varphi)}{\sin^2\alpha_1}\times\\
&\times[f^1(\varphi_1,\varphi_1^0)e^{-ika\sqrt{\lambda^2+\mu^2}}-if^1(\varphi_2,\varphi_2^0)e^{ika\sqrt{\lambda^2+\mu^2}}]\biggr\},\\
E_\vartheta=H_\varphi={}&\frac{ae^{i\pi/4}}{\sqrt{2\pi ka\sqrt{\lambda^2+\mu^2}}}\frac{e^{ikR}}R\times\\
&\times\biggl\{-H_{0x_1}\frac{\cos(\psi_1-\varphi)}{\sin^2\alpha_1}[g^1(\varphi_1,\varphi_1^0)e^{-ika\sqrt{\lambda^2+\mu^2}}\\
&-ig^1(\varphi_2,\varphi_2^0)e^{ika\sqrt{\lambda^2+\mu^2}}]+E_{0x_1}\cos\vartheta\frac{\sin(\psi_1-\varphi)}{\sin^2\alpha_1}\times\\
&\times[f^1(\varphi_1,\varphi_1^0)e^{-ika\sqrt{\lambda^2+\mu^2}}-if^1(\varphi_2,\varphi_2^0)e^{ika\sqrt{\lambda^2+\mu^2}}]\biggr\},
\end{aligned}\tag{11.32}
\]
if we use the identities
\[
\left.\begin{aligned}
\frac{\sin(\psi_1-\varphi)}{\sin\alpha_1\cos\varphi_1}-\frac{\cos(\psi_1-\varphi)}{\sin^2\alpha_1}\cos\alpha_1\operatorname{tg}\varphi_1&\equiv\frac{\cos\vartheta\sin(\psi_1-\varphi)}{\sin^2\alpha_1},\\
\frac{\cos(\psi_1-\varphi)}{\sin\alpha_1\cos\varphi_1}+\frac{\sin(\psi_1-\varphi)}{\sin^2\alpha_1}\cos\alpha_1\operatorname{tg}\varphi_1&\equiv\frac{\cos(\psi_1-\varphi)}{\sin^2\alpha_1\cos\vartheta}.
\end{aligned}\right\}\tag{11.33}
\]
The operations carried out above may be briefly summarized in the following way. The field from the nonuniform part of the current on the disk
\[\mathbf A=\frac ac\frac{e^{ikR}}R\int_0^{2\pi}d\psi\int_0^a\mathbf J(\rho,\psi)e^{ik\rho\Phi}\,d\rho\]
is found (without direct calculation of the current) in terms of the known field of an auxiliary half-plane
\end{scanpage}
